% =============================================================================
% Related Work section of "Portuguese-MMLU: Measuring Massive Multitask
% Language Understanding in Portuguese".
%
% Usage: % =============================================================================
% Related Work section of "Portuguese-MMLU: Measuring Massive Multitask
% Language Understanding in Portuguese".
%
% Usage: % =============================================================================
% Related Work section of "Portuguese-MMLU: Measuring Massive Multitask
% Language Understanding in Portuguese".
%
% Usage: % =============================================================================
% Related Work section of "Portuguese-MMLU: Measuring Massive Multitask
% Language Understanding in Portuguese".
%
% Usage: \input{related_work} in acl_latex.tex, after the Introduction.
% Required packages: natbib (\citep, \citet; loaded by acl.sty) and hyperref
% or url (\url in footnotes; the ACL template loads hyperref).
% Bibliography entries are in related_work.bib: append them to custom.bib or
% use \bibliography{custom,related_work}. Every cited work is peer-reviewed;
% web resources appear only in footnotes.
% =============================================================================

\section{Related Work}
\label{sec:related}

MMLU \citep{hendrycks2021mmlu} established the format that most knowledge
benchmarks follow today: multiple-choice questions over 57 academic and
professional subjects, collected from exams and study materials and scored by
accuracy in zero- and few-shot settings. Its limitations have driven two lines
of follow-up work. The first revises the English benchmark itself: MMLU-Pro
raises the number of options from four to ten and removes trivial and noisy
items, lowering the accuracy of frontier models by 16 to 33 points
\citep{wang2024mmlupro}, and MMLU-Redux re-annotates 5,700 questions by hand,
estimates that 6.5\% of MMLU items are erroneous and shows that correcting
them changes model rankings \citep{gema2025mmluredux}. The second builds
benchmarks directly from official human examinations, whose items are written
and calibrated by experts for a population of real candidates: AGIEval, for
instance, gathers college-entrance exams, law-school admission tests, bar
exams and mathematics competitions in English and Chinese, and finds that
models remain weak on questions that demand domain knowledge or multi-step
reasoning \citep{zhong2024agieval}. Portuguese-MMLU belongs to this second
line, since all its questions come from official Brazilian exams, and the two
revision efforts above motivate the quality controls applied before release.

The usual way to evaluate a language other than English is to translate an
English benchmark. Global-MMLU translates MMLU into 42 languages, Portuguese
among them, with professional and community translators and annotators, and
labels 28\% of its questions as culturally sensitive; model rankings change
when only that subset is used, which the authors attribute to the
Western-centric content of the source \citep{singh2025globalmmlu}. MMLU-ProX
extends MMLU-Pro to 29 languages, including Portuguese, through model
translation followed by expert review, and reports gaps of up to 24 points
between high- and low-resource languages \citep{xuan2025mmluprox}; OpenAI's
MMMLU offers a professionally translated Brazilian Portuguese version of
MMLU.\footnote{\url{https://huggingface.co/datasets/openai/MMMLU}} Native
benchmarks take the opposite route and argue that translation imports errors
and a foreign curriculum. CMMLU \citep{li2024cmmlu}, ArabicMMLU
\citep{koto2024arabicmmlu}, TurkishMMLU \citep{yuksel2024turkishmmlu} and
KMMLU \citep{son2025kmmlu} are built from exams and teaching materials written
in the target language; KMMLU, with 35,030 questions from Korean examinations,
finds that models score below the pass mark of the source exams and that one
fifth of its questions require knowledge of Korean culture. Portuguese appears
in multi-country exam collections only as a small slice: M3Exam includes about
1,400 Brazilian questions, among them ENEM items, in a nine-language benchmark
\citep{zhang2023m3exam}; INCLUDE has 581 Portuguese questions among 197,243
across 44 languages \citep{romanou2025include}; and Kaleidoscope, a
multimodal collection in 18 languages, contributes 2,000 Portuguese items,
half of which require an image \citep{salazar2026kaleidoscope}. A native,
multi-domain benchmark of the scale of CMMLU or KMMLU does not exist for
Portuguese.

Brazilian exams were proposed as a test for artificial intelligence before
large language models: the ENEM Challenge turned the national high-school
exam into a machine-readable question set and reported retrieval and
word-embedding solvers between 26\% and 29\% accuracy on its text-only
questions \citep{silveira2018enem}. With generative models, the exams became
the core of Portuguese evaluation suites. BLUEX collects the entrance exams of
USP and UNICAMP \citep{almeida2023bluex}; the Poeta suite, with 14 Portuguese
datasets, was used to show that continued pretraining in Portuguese improves
the Sabi\'a models over multilingual baselines \citep{pires2023sabia}; and the
Tucano models were evaluated with a harness that combines ENEM, BLUEX and the
OAB bar exam with inference, similarity and classification tasks and four
benchmarks machine-translated from English \citep{correa2025tucano}.
Recent resources widen the formats and the variants covered: EduBench
contains 3,149 discursive questions from entrance exams with expert reference
answers \citep{paiola2026edubench}; a benchmark built from the 2025 ENAMED
medical-licensing exam scores 22 models with the exam's official
item-response model and finds that many open-weight models fall below the
passing threshold \citep{correia2026enamed}; and CLARIN-PT-LDB is a
leaderboard for European Portuguese \citep{silva2026clarinptldb}. These
resources are either single-exam, with tens to a few thousand questions, or
built around the same three sources, ENEM, BLUEX and OAB, so that subjects
such as health, business, engineering and the social sciences at
undergraduate level remain largely untested in Portuguese.

Portuguese-MMLU differs from this prior work in scope and construction. It
gathers native questions from 17 Brazilian examinations, spanning university
entrance exams, an informatics olympiad, professional certification,
public-sector recruitment and medical licensing and residency exams, at two
academic levels. Every question is assigned a subject and a macro area,
duplicates are removed at the exact and the semantic level, and items whose
statement cannot be recovered are revised out before release, so that results
can be compared per area across the many disciplines the exams cover.
 in acl_latex.tex, after the Introduction.
% Required packages: natbib (\citep, \citet; loaded by acl.sty) and hyperref
% or url (\url in footnotes; the ACL template loads hyperref).
% Bibliography entries are in related_work.bib: append them to custom.bib or
% use \bibliography{custom,related_work}. Every cited work is peer-reviewed;
% web resources appear only in footnotes.
% =============================================================================

\section{Related Work}
\label{sec:related}

MMLU \citep{hendrycks2021mmlu} established the format that most knowledge
benchmarks follow today: multiple-choice questions over 57 academic and
professional subjects, collected from exams and study materials and scored by
accuracy in zero- and few-shot settings. Its limitations have driven two lines
of follow-up work. The first revises the English benchmark itself: MMLU-Pro
raises the number of options from four to ten and removes trivial and noisy
items, lowering the accuracy of frontier models by 16 to 33 points
\citep{wang2024mmlupro}, and MMLU-Redux re-annotates 5,700 questions by hand,
estimates that 6.5\% of MMLU items are erroneous and shows that correcting
them changes model rankings \citep{gema2025mmluredux}. The second builds
benchmarks directly from official human examinations, whose items are written
and calibrated by experts for a population of real candidates: AGIEval, for
instance, gathers college-entrance exams, law-school admission tests, bar
exams and mathematics competitions in English and Chinese, and finds that
models remain weak on questions that demand domain knowledge or multi-step
reasoning \citep{zhong2024agieval}. Portuguese-MMLU belongs to this second
line, since all its questions come from official Brazilian exams, and the two
revision efforts above motivate the quality controls applied before release.

The usual way to evaluate a language other than English is to translate an
English benchmark. Global-MMLU translates MMLU into 42 languages, Portuguese
among them, with professional and community translators and annotators, and
labels 28\% of its questions as culturally sensitive; model rankings change
when only that subset is used, which the authors attribute to the
Western-centric content of the source \citep{singh2025globalmmlu}. MMLU-ProX
extends MMLU-Pro to 29 languages, including Portuguese, through model
translation followed by expert review, and reports gaps of up to 24 points
between high- and low-resource languages \citep{xuan2025mmluprox}; OpenAI's
MMMLU offers a professionally translated Brazilian Portuguese version of
MMLU.\footnote{\url{https://huggingface.co/datasets/openai/MMMLU}} Native
benchmarks take the opposite route and argue that translation imports errors
and a foreign curriculum. CMMLU \citep{li2024cmmlu}, ArabicMMLU
\citep{koto2024arabicmmlu}, TurkishMMLU \citep{yuksel2024turkishmmlu} and
KMMLU \citep{son2025kmmlu} are built from exams and teaching materials written
in the target language; KMMLU, with 35,030 questions from Korean examinations,
finds that models score below the pass mark of the source exams and that one
fifth of its questions require knowledge of Korean culture. Portuguese appears
in multi-country exam collections only as a small slice: M3Exam includes about
1,400 Brazilian questions, among them ENEM items, in a nine-language benchmark
\citep{zhang2023m3exam}; INCLUDE has 581 Portuguese questions among 197,243
across 44 languages \citep{romanou2025include}; and Kaleidoscope, a
multimodal collection in 18 languages, contributes 2,000 Portuguese items,
half of which require an image \citep{salazar2026kaleidoscope}. A native,
multi-domain benchmark of the scale of CMMLU or KMMLU does not exist for
Portuguese.

Brazilian exams were proposed as a test for artificial intelligence before
large language models: the ENEM Challenge turned the national high-school
exam into a machine-readable question set and reported retrieval and
word-embedding solvers between 26\% and 29\% accuracy on its text-only
questions \citep{silveira2018enem}. With generative models, the exams became
the core of Portuguese evaluation suites. BLUEX collects the entrance exams of
USP and UNICAMP \citep{almeida2023bluex}; the Poeta suite, with 14 Portuguese
datasets, was used to show that continued pretraining in Portuguese improves
the Sabi\'a models over multilingual baselines \citep{pires2023sabia}; and the
Tucano models were evaluated with a harness that combines ENEM, BLUEX and the
OAB bar exam with inference, similarity and classification tasks and four
benchmarks machine-translated from English \citep{correa2025tucano}.
Recent resources widen the formats and the variants covered: EduBench
contains 3,149 discursive questions from entrance exams with expert reference
answers \citep{paiola2026edubench}; a benchmark built from the 2025 ENAMED
medical-licensing exam scores 22 models with the exam's official
item-response model and finds that many open-weight models fall below the
passing threshold \citep{correia2026enamed}; and CLARIN-PT-LDB is a
leaderboard for European Portuguese \citep{silva2026clarinptldb}. These
resources are either single-exam, with tens to a few thousand questions, or
built around the same three sources, ENEM, BLUEX and OAB, so that subjects
such as health, business, engineering and the social sciences at
undergraduate level remain largely untested in Portuguese.

Portuguese-MMLU differs from this prior work in scope and construction. It
gathers native questions from 17 Brazilian examinations, spanning university
entrance exams, an informatics olympiad, professional certification,
public-sector recruitment and medical licensing and residency exams, at two
academic levels. Every question is assigned a subject and a macro area,
duplicates are removed at the exact and the semantic level, and items whose
statement cannot be recovered are revised out before release, so that results
can be compared per area across the many disciplines the exams cover.
 in acl_latex.tex, after the Introduction.
% Required packages: natbib (\citep, \citet; loaded by acl.sty) and hyperref
% or url (\url in footnotes; the ACL template loads hyperref).
% Bibliography entries are in related_work.bib: append them to custom.bib or
% use \bibliography{custom,related_work}. Every cited work is peer-reviewed;
% web resources appear only in footnotes.
% =============================================================================

\section{Related Work}
\label{sec:related}

MMLU \citep{hendrycks2021mmlu} established the format that most knowledge
benchmarks follow today: multiple-choice questions over 57 academic and
professional subjects, collected from exams and study materials and scored by
accuracy in zero- and few-shot settings. Its limitations have driven two lines
of follow-up work. The first revises the English benchmark itself: MMLU-Pro
raises the number of options from four to ten and removes trivial and noisy
items, lowering the accuracy of frontier models by 16 to 33 points
\citep{wang2024mmlupro}, and MMLU-Redux re-annotates 5,700 questions by hand,
estimates that 6.5\% of MMLU items are erroneous and shows that correcting
them changes model rankings \citep{gema2025mmluredux}. The second builds
benchmarks directly from official human examinations, whose items are written
and calibrated by experts for a population of real candidates: AGIEval, for
instance, gathers college-entrance exams, law-school admission tests, bar
exams and mathematics competitions in English and Chinese, and finds that
models remain weak on questions that demand domain knowledge or multi-step
reasoning \citep{zhong2024agieval}. Portuguese-MMLU belongs to this second
line, since all its questions come from official Brazilian exams, and the two
revision efforts above motivate the quality controls applied before release.

The usual way to evaluate a language other than English is to translate an
English benchmark. Global-MMLU translates MMLU into 42 languages, Portuguese
among them, with professional and community translators and annotators, and
labels 28\% of its questions as culturally sensitive; model rankings change
when only that subset is used, which the authors attribute to the
Western-centric content of the source \citep{singh2025globalmmlu}. MMLU-ProX
extends MMLU-Pro to 29 languages, including Portuguese, through model
translation followed by expert review, and reports gaps of up to 24 points
between high- and low-resource languages \citep{xuan2025mmluprox}; OpenAI's
MMMLU offers a professionally translated Brazilian Portuguese version of
MMLU.\footnote{\url{https://huggingface.co/datasets/openai/MMMLU}} Native
benchmarks take the opposite route and argue that translation imports errors
and a foreign curriculum. CMMLU \citep{li2024cmmlu}, ArabicMMLU
\citep{koto2024arabicmmlu}, TurkishMMLU \citep{yuksel2024turkishmmlu} and
KMMLU \citep{son2025kmmlu} are built from exams and teaching materials written
in the target language; KMMLU, with 35,030 questions from Korean examinations,
finds that models score below the pass mark of the source exams and that one
fifth of its questions require knowledge of Korean culture. Portuguese appears
in multi-country exam collections only as a small slice: M3Exam includes about
1,400 Brazilian questions, among them ENEM items, in a nine-language benchmark
\citep{zhang2023m3exam}; INCLUDE has 581 Portuguese questions among 197,243
across 44 languages \citep{romanou2025include}; and Kaleidoscope, a
multimodal collection in 18 languages, contributes 2,000 Portuguese items,
half of which require an image \citep{salazar2026kaleidoscope}. A native,
multi-domain benchmark of the scale of CMMLU or KMMLU does not exist for
Portuguese.

Brazilian exams were proposed as a test for artificial intelligence before
large language models: the ENEM Challenge turned the national high-school
exam into a machine-readable question set and reported retrieval and
word-embedding solvers between 26\% and 29\% accuracy on its text-only
questions \citep{silveira2018enem}. With generative models, the exams became
the core of Portuguese evaluation suites. BLUEX collects the entrance exams of
USP and UNICAMP \citep{almeida2023bluex}; the Poeta suite, with 14 Portuguese
datasets, was used to show that continued pretraining in Portuguese improves
the Sabi\'a models over multilingual baselines \citep{pires2023sabia}; and the
Tucano models were evaluated with a harness that combines ENEM, BLUEX and the
OAB bar exam with inference, similarity and classification tasks and four
benchmarks machine-translated from English \citep{correa2025tucano}.
Recent resources widen the formats and the variants covered: EduBench
contains 3,149 discursive questions from entrance exams with expert reference
answers \citep{paiola2026edubench}; a benchmark built from the 2025 ENAMED
medical-licensing exam scores 22 models with the exam's official
item-response model and finds that many open-weight models fall below the
passing threshold \citep{correia2026enamed}; and CLARIN-PT-LDB is a
leaderboard for European Portuguese \citep{silva2026clarinptldb}. These
resources are either single-exam, with tens to a few thousand questions, or
built around the same three sources, ENEM, BLUEX and OAB, so that subjects
such as health, business, engineering and the social sciences at
undergraduate level remain largely untested in Portuguese.

Portuguese-MMLU differs from this prior work in scope and construction. It
gathers native questions from 17 Brazilian examinations, spanning university
entrance exams, an informatics olympiad, professional certification,
public-sector recruitment and medical licensing and residency exams, at two
academic levels. Every question is assigned a subject and a macro area,
duplicates are removed at the exact and the semantic level, and items whose
statement cannot be recovered are revised out before release, so that results
can be compared per area across the many disciplines the exams cover.
 in acl_latex.tex, after the Introduction.
% Required packages: natbib (\citep, \citet; loaded by acl.sty) and hyperref
% or url (\url in footnotes; the ACL template loads hyperref).
% Bibliography entries are in related_work.bib: append them to custom.bib or
% use \bibliography{custom,related_work}. Every cited work is peer-reviewed;
% web resources appear only in footnotes.
% =============================================================================

\section{Related Work}
\label{sec:related}

MMLU \citep{hendrycks2021mmlu} established the format that most knowledge
benchmarks follow today: multiple-choice questions over 57 academic and
professional subjects, collected from exams and study materials and scored by
accuracy in zero- and few-shot settings. Its limitations have driven two lines
of follow-up work. The first revises the English benchmark itself: MMLU-Pro
raises the number of options from four to ten and removes trivial and noisy
items, lowering the accuracy of frontier models by 16 to 33 points
\citep{wang2024mmlupro}, and MMLU-Redux re-annotates 5,700 questions by hand,
estimates that 6.5\% of MMLU items are erroneous and shows that correcting
them changes model rankings \citep{gema2025mmluredux}. The second builds
benchmarks directly from official human examinations, whose items are written
and calibrated by experts for a population of real candidates: AGIEval, for
instance, gathers college-entrance exams, law-school admission tests, bar
exams and mathematics competitions in English and Chinese, and finds that
models remain weak on questions that demand domain knowledge or multi-step
reasoning \citep{zhong2024agieval}. Portuguese-MMLU belongs to this second
line, since all its questions come from official Brazilian exams, and the two
revision efforts above motivate the quality controls applied before release.

The usual way to evaluate a language other than English is to translate an
English benchmark. Global-MMLU translates MMLU into 42 languages, Portuguese
among them, with professional and community translators and annotators, and
labels 28\% of its questions as culturally sensitive; model rankings change
when only that subset is used, which the authors attribute to the
Western-centric content of the source \citep{singh2025globalmmlu}. MMLU-ProX
extends MMLU-Pro to 29 languages, including Portuguese, through model
translation followed by expert review, and reports gaps of up to 24 points
between high- and low-resource languages \citep{xuan2025mmluprox}; OpenAI's
MMMLU offers a professionally translated Brazilian Portuguese version of
MMLU.\footnote{\url{https://huggingface.co/datasets/openai/MMMLU}} Native
benchmarks take the opposite route and argue that translation imports errors
and a foreign curriculum. CMMLU \citep{li2024cmmlu}, ArabicMMLU
\citep{koto2024arabicmmlu}, TurkishMMLU \citep{yuksel2024turkishmmlu} and
KMMLU \citep{son2025kmmlu} are built from exams and teaching materials written
in the target language; KMMLU, with 35,030 questions from Korean examinations,
finds that models score below the pass mark of the source exams and that one
fifth of its questions require knowledge of Korean culture. Portuguese appears
in multi-country exam collections only as a small slice: M3Exam includes about
1,400 Brazilian questions, among them ENEM items, in a nine-language benchmark
\citep{zhang2023m3exam}; INCLUDE has 581 Portuguese questions among 197,243
across 44 languages \citep{romanou2025include}; and Kaleidoscope, a
multimodal collection in 18 languages, contributes 2,000 Portuguese items,
half of which require an image \citep{salazar2026kaleidoscope}. A native,
multi-domain benchmark of the scale of CMMLU or KMMLU does not exist for
Portuguese.

Brazilian exams were proposed as a test for artificial intelligence before
large language models: the ENEM Challenge turned the national high-school
exam into a machine-readable question set and reported retrieval and
word-embedding solvers between 26\% and 29\% accuracy on its text-only
questions \citep{silveira2018enem}. With generative models, the exams became
the core of Portuguese evaluation suites. BLUEX collects the entrance exams of
USP and UNICAMP \citep{almeida2023bluex}; the Poeta suite, with 14 Portuguese
datasets, was used to show that continued pretraining in Portuguese improves
the Sabi\'a models over multilingual baselines \citep{pires2023sabia}; and the
Tucano models were evaluated with a harness that combines ENEM, BLUEX and the
OAB bar exam with inference, similarity and classification tasks and four
benchmarks machine-translated from English \citep{correa2025tucano}.
Recent resources widen the formats and the variants covered: EduBench
contains 3,149 discursive questions from entrance exams with expert reference
answers \citep{paiola2026edubench}; a benchmark built from the 2025 ENAMED
medical-licensing exam scores 22 models with the exam's official
item-response model and finds that many open-weight models fall below the
passing threshold \citep{correia2026enamed}; and CLARIN-PT-LDB is a
leaderboard for European Portuguese \citep{silva2026clarinptldb}. These
resources are either single-exam, with tens to a few thousand questions, or
built around the same three sources, ENEM, BLUEX and OAB, so that subjects
such as health, business, engineering and the social sciences at
undergraduate level remain largely untested in Portuguese.

Portuguese-MMLU differs from this prior work in scope and construction. It
gathers native questions from 17 Brazilian examinations, spanning university
entrance exams, an informatics olympiad, professional certification,
public-sector recruitment and medical licensing and residency exams, at two
academic levels. Every question is assigned a subject and a macro area,
duplicates are removed at the exact and the semantic level, and items whose
statement cannot be recovered are revised out before release, so that results
can be compared per area across the many disciplines the exams cover.
